\section{The cyclotomic coordinate of the finite gentle levels}
\label{sec:gt-certificates}

Write $U_N=(\mathbb Z/N\mathbb Z)^\times$, and use natural logarithms.
The word ``level'' in the supplied material can mean a matrix modulus,
a power imposed on a relative kernel, the order of a finite group, or
the modulus of the multiplier itself. These are different data. We
retain the actual quotient maps when making the comparison.

\subsection{The precise Bach input}

\paragraph{Finding T01 (the corrected uniform certificate).}
Under Dirichlet GRH, and hence under Alp\"oge's theorem as stipulated
in this problem, for every integer $N\geq3$ the set
\[
 \mathcal P_3(N)=\{p:p\text{ prime},\ p\nmid N,
                              \ p<3(\log N)^2\}
 \tag{T01}
\]
generates $U_N$. The groups $U_1,U_2$ are trivial and need no generators.
The inequality is strict, and there are no omitted small moduli.
Neither a Hecke GRH over a nonabelian field nor Masullo's claim is needed.

The coefficient $2$ attributed to unit generation in the prompt and
Fable's suggestion~7 is not the statement of Bach's cited theorem.
Bach's Theorem~2, printed page~372, treats a proper subgroup containing
every positive integer below a threshold; integers not coprime to the
modulus are allowed to be outside the subgroup. Its bound is
$2(\log N)^2$. Theorem~3 on the same page imposes coprimality and has
bound $3(\log N)^2$. The consequence for generators is stated on
printed page~373 \cite[Theorems~2--3, pp.~372--373]{Bach1990}.
The distinction is also explicitly recorded in
\cite[\S1.1]{LLS2015}.

\begin{proof}
Only the passage from Bach's integer statement to prime generators
requires proof. Let $H$ be generated by $\mathcal P_3(N)$. Every positive
integer $a<3(\log N)^2$ coprime to $N$ factors into primes in
$\mathcal P_3(N)$, so its residue lies in $H$. If $H$ were proper,
Bach's Theorem~3, applied with $x=3(\log N)^2$, would give
$x<3(\log N)^2$, a contradiction. This proves the assertion. For
$N=1,2$ the assertion is the direct trivial-group calculation.

The distinction between the two theorems cannot be discarded: in
$U_{12}$ the subgroup $\{1,5\}$ omits the positive integer $2$, but
$2$ is not a unit; the least omitted positive unit is $7$. This example
refutes the inference that a witness supplied by Theorem~2 must be a
unit. It does not refute the stronger numerical coefficient-$2$
generation statement, which is not asserted false here.
\end{proof}

The cited theorem's analytic proof has not been independently reproved
in this run, and Bach's original computation through one million moduli
has not been rerun. Its exact hypotheses, constants, strict inequalities,
and small-modulus scope were checked in the original scanned article.
The direct computations below are independent of this analytic input.

\paragraph{Finding T07 (a sharper criterion and a large-modulus consequence).}
Theorem~1.1 of Lamzouri--Li--Soundararajan, in the authors' corrected
arXiv version~v3 of 8~December~2016, applies to every proper subgroup
of $U_N$, for $N\geq3000$, with the same Dirichlet GRH input.
Define
\[
 \begin{split}
 A(N)&=\max\left(0,2\log\log N-\frac85
                    -\sum_{p\mid N}\frac{\log p}{p-1}\right),\\
 B(N)&=\max\left(0,2\log\log N+3
          +\frac{2\omega(N)(\log\log N)^2}{\log N}-2A(N)\right),
 \end{split}
\]
where $\omega(N)$ counts distinct prime divisors. A prime $\ell\nmid N$
outside the subgroup exists with
$\ell\leq(\log N+B(N))^2$
\cite[Theorem~1.1(1), corrected arXiv version~v3]{LLS2015}.
Thus the desired strict coefficient-$2$ generator bound follows
whenever $B(N)<(\sqrt2-1)\log N$, by the same subgroup argument as
T01. Part~(2) gives $\ell\leq(\log N)^2$ under the additional
condition that $N$ has no prime divisor below $(\log N)^2$.
Theorem~1.4 instead concerns a specified coset and depends on its
subgroup index; Corollary~1.1 concerns quadratic nonresidues for prime
moduli. Neither states a uniform coefficient-$2$ generator bound for
arbitrary composite moduli.

The displayed bound in Theorem~1.1(1) is not uniformly below
$2(\log N)^2$ merely for $N\geq3000$. Put
$R(N)=(\log N+B(N))^2/(\log N)^2$. Exact rational interval calculations
give
\[
 \begin{split}
 4.98743&<R(30030)<4.98744,\\
 2.02303&<R(13082761331670030)<2.02304.
 \end{split}
\]
The moduli are the products of the first six and first fourteen primes,
respectively. The executable \texttt{check\_lls\_bound\_scope.py}
certifies these intervals using the positive logarithm series of
(T05.1), also applied to rational endpoints for $\log\log N$;
its two records are in \texttt{LLS\_BOUND\_EXPRESSION\_CHECKS.tsv}.
These are evaluations of the theorem's upper-bound expression,
\emph{not} counterexamples to actual coefficient-$2$ generation.
They show why the displayed LLS theorem plus checking only $N<3000$
does not complete such a uniform claim. The authors' v3 revision notice
identifies corrections to the asymptotic Theorems~1.2--1.3; no constants
from their superseded version are used here.

Under A, the strict coefficient-$2$ generators also suffice for every
integer $N\ge\exp(300)$. This deliberately coarse threshold is a new
deduction of the displayed LLS bound, not a threshold quoted from that
paper; the remaining interval is not asserted checked.

\begin{proof}
The sufficient $B(N)$ criterion follows by the proper-subgroup argument
in T01. The two expression enclosures follow from the exact rational
positive-series estimates and complete calculations just specified.
For the final assertion put $L=\log N\ge300$ and $l=\log L$.
Since $2^{\omega(N)}\le N$ and $\log2>2/3$, the defining formula gives
$B(N)<3l^2+2l+3$. Indeed $A(N)\ge0$ and
$2\omega(N)l^2/L<3l^2$. The function
$F(L)=(3(\log L)^2+2\log L+3)/L$ has derivative
$-(3\log L-1)(\log L-1)/L^2<0$ for $L\ge300$.
Also $\log300<6$, since $e>8/3$ and $(8/3)^6>300$.
Thus $F(L)\le F(300)<123/300=41/100$ and
\[
 (L+B(N))^2<(141/100)^2L^2
             =\frac{19881}{10000}L^2<2L^2.
\]
LLS Theorem~1.1 applies since $N\ge\exp(300)>3000$ and gives a
coprime prime outside any proper subgroup below this last bound.
The same subgroup argument proves generation. The elementary lower
bound $\log2>2/3$ is, for example, the first term plus the positive
remainder in $\log2=2\sum_{j\ge0}(1/3)^{2j+1}/(2j+1)$.
\end{proof}

\subsection{What the finite certificate certifies}

\paragraph{Finding T02 (finite cyclotomic tests and their limits).}
Let $\mathcal G=\widehat{\mathrm{GT}}_0$, with multiplier
$\lambda:\mathcal G\to\widehat{\mathbb Z}^{\times}$, and let
$q:\mathcal G\twoheadrightarrow Q$ be a finite quotient. To speak of
its full multiplier modulo $N$, require an actual homomorphism
$\lambda_N:Q\to U_N$ with
\[
                    \lambda\bmod N=\lambda_N\circ q.
 \tag{T02.1}
\]
Equivalently, $\ker q\subseteq\ker(\lambda\bmod N)$. If this is not
available, the actual finite image of $(q,\lambda\bmod N)$ is a
specified refinement that does retain the multiplier; this must not
silently be identified with $Q$.

Suppose $P$ is any proved generating set of $U_N$, for instance
$\mathcal P_3(N)$ under T01, or a list certified in T05. Then:
\begin{enumerate}
\item a subgroup $J\leq Q$ has $\lambda_N(J)=U_N$ if its multiplier
image contains every residue in $P$;
\item two homomorphisms from $U_N$ agree if they agree on $P$;
\item a nontrivial character of $U_N$ is nontrivial on some member of $P$.
\end{enumerate}
These statements concern subgroups and homomorphisms already shown to
have the indicated algebraic properties. A set of candidate shadows
has not thereby been shown to be a subgroup. The certificate does not
check the two cycle equations, the pentagon, the $f$-coordinate, or
compatible lifting through all finite quotients.

\begin{proof}
Each assertion follows by writing every unit as a word in $P$. For
the third, otherwise the kernel of the character would contain all of
$U_N$. For the first, it is essential that the image be a subgroup,
so it contains all words in its specified generators. The factorization
criterion in (T02.1) is the elementary universal property of a quotient.

For the limitation, retain the supplied LM7\_PUSH E.65 theorem:
there is a continuous character $\Delta:\mathcal G\to\mathbb F_2$
vanishing on the Galois subgroup and taking value $1$ on an actual
element of $\ker\lambda$ \cite[E.58, E.65]{LM7run}. It cannot factor
through $\lambda$ or through any $\lambda\bmod N$, since a homomorphism
of that form is zero on $\ker\lambda$. Moreover $\ker\Delta$ is proper
but has the same full cyclotomic image as $\mathcal G$: it contains
$G_{\mathbb Q}$, whose cyclotomic character is onto. Thus even all of
the cyclotomic finite certificates together cannot distinguish these
two subgroups. This is a concrete counterexample within the supplied
gentle-group framework to inferring equality of groups from equality
of their multiplier images. It uses the supplied E.65 theorem as
background, not a new proof of E.58's profinite construction.
\end{proof}

For clarity, the arithmetic inclusion already gives
$\lambda(\mathcal G)=\widehat{\mathbb Z}^{\times}$, independently of
GRH. When an actual finite quotient retains $\lambda\bmod N$, its
full multiplier image is therefore already $U_N$. Bach's input gives
a bounded list of witnesses or tests, not the original surjectivity
theorem. Also, an arithmetic Frobenius at $p\nmid N$ has multiplier
$p$ modulo $N$; no claim is made that its full profinite multiplier
is the nonunit integer $p$ in $\widehat{\mathbb Z}$.

\paragraph{The scheme and fibre-functor formulation.}
Put $S_N=\operatorname{Spec}\mathbb Z[1/N]$ and
$X_N=\operatorname{Spec}\mathbb Z[1/N,\zeta_N]$. This is the connected
finite \'etale Galois cover with group $U_N$. Indeed it is the primitive
root component of the \'etale scheme of roots of $T^N-1$ after $N$ is
inverted; the cyclotomic field identifies its generic Galois group.
Arithmetic Frobenius sends $\zeta_N$ to $\zeta_N^p$. T02 is consequently
a certificate on a finite quotient of the fibre-functor automorphism
group of $S_N$, in SGA~1, Expos\'e~V, Th\'eor\`eme~4.1 and \S7
\cite{SGA1}. This is the applicable Grothendieck framework. It supplies
no additional purity, crystalline, divided-power, or motivic assertion.
For each prime $\ell$ the compatible roots of unity are the Tate module
$\mathbb Z_\ell(1)$; finite level records its cyclotomic scalar action.

\subsection{Elementary certificates for the two- and three-primary towers}

\paragraph{Finding T03 (strict coefficient-$2$ certificates without GRH).}
For integers $a,b\geq0$ and $N=2^a3^b>1$, the following prime lists generate $U_N$, and every
listed prime is coprime to $N$ and strictly less than $2(\log N)^2$:
\[
\begin{array}{c|c}
 (a,b)&\text{prime generators}\\\hline
 (1,0)&\varnothing\\
 (2,0)&3\\
 a\geq3,\ b=0&3,5\\
 a=0,\ b\geq1&2\\
 a=1,\ b\geq1&5\\
 a=2,\ b\geq1&5,7\\
 a\geq3,\ b\geq1&5,7,13
\end{array}
 \tag{T03}
\]
Also $2$ generates $U_{5^b}$ for every $b\geq1$ and satisfies the same
strict bound. These are elementary unconditional statements, requiring
neither of the analytic inputs in the prompt.

\begin{proof}
For $a\geq3$, the subgroup of units congruent to $1$ modulo $4$ has
order $2^{a-2}$ and is generated by $5$. The elementary valuation
identity
\[
                  v_2(5^{2^j}-1)=j+2
 \tag{T03.1}
\]
proves the order: the base case is $5-1=4$, and on each squaring the
factor $5^{2^j}+1$ has valuation one. Since $3\equiv-1\pmod4$, the
two primes $3,5$ generate all of $U_{2^a}$. The case $a=2$ is direct.

Both $2$ and $5$ have order $2\cdot3^{b-1}$ modulo $3^b$. To see this,
they are $-1$ modulo three, and
\[
 v_3(2^2-1)=v_3(5^2-1)=1,
 \qquad v_3(z^3-1)=v_3(z-1)+1\quad(z\equiv1\pmod3).
 \tag{T03.2}
\]
The latter identity follows by expanding $(1+3^s u)^3-1$ for
$s\geq1$ and $3\nmid u$: its first term has uniquely least
valuation $s+1$. Repeated use gives the asserted orders, equal to
$|U_{3^b}|$. This proves the first two ternary cases.

For $a=2,b\geq1$, the quotient of $U_N$ by its squares is canonically
$U_{12}$, generated by $5,7$: their products give $1,5,7,11$.
The projection of $5$ generates the ternary primary component, by (T03.2).
Thus the subgroup they generate is all of $U_N$. Here is a proof
that these two checks suffice: if $H$ is proper in a finite abelian
group $G$, then $G/H$ has a quotient of prime order; the only
possible primes dividing $|U_N|$ are two and three. Such a quotient
would contradict one of the two checks.

For $a\geq3,b\geq1$, the quotient by squares is instead $U_{24}$.
Its eight distinct elements are the products of $5,7,13$:
\[
                    1,5,7,13,11,17,19,23.
 \tag{T03.3}
\]
Reduction from $U_{2^a}$ to $U_8$ kills exactly its squares, using
(T03.1); reduction from $U_{3^b}$ to $U_3$ does the same for its
two-primary quotient, using (T03.2). This proves the assertion about
the quotient by squares. The projection of $5$ again generates the
three-primary component. The preceding prime-order-quotient argument
finishes the mixed case. Empty or absent primary components cause
no exception.

For $5^b$, $2$ has order four modulo five, and $v_5(2^4-1)=1$.
The same binomial argument, now with fifth powers, shows
$v_5(2^{4\cdot5^j}-1)=j+1$, so its order is $4\cdot5^{b-1}$.

It remains to verify the strict bounds for these entire families,
rather than just for a finite sample. The elementary series
$\log2=2(1/3+(1/3)^3/3+\cdots)$ gives $\log2>2/3$; also
$e<3$, by its factorial series, gives $\log3>1$. For the minimal
moduli $3,4,8,6,12,24$ in the nonempty rows, respectively,
these inequalities give
\[
 2(\log N)^2>2,\quad 32/9,\quad8,\quad50/9,\quad98/9,\quad18.
\]
These exceed respectively the largest proposed primes $2,3,5,5,7,13$.
The bound increases with $N$. For $5^b$ it suffices that $5>3$.
This proves strictness in every case, including all boundary cases.
\end{proof}

\subsection{Identifying the actual scalar moduli}

For the LM7 homology tower, retain the actual maps
\[
 B=\langle s,r\mid s^2=r^3=1\rangle,
 \quad t=sr,\quad
 q:B\longrightarrow\operatorname{SL}_2(\mathbb Z/12)/\{\pm I\},
\]
where $q(s)$ and $q(r)$ are the classes of
$S=\left(\begin{smallmatrix}0&1\\-1&0\end{smallmatrix}\right)$ and
$R=\left(\begin{smallmatrix}0&-1\\1&-1\end{smallmatrix}\right)$.
Set $N_B=\ker q$, $C=N_B^{\rm ab}$, and
$G_j=B/(N_B'N_B^j)$. These are the discrete source definitions of
LM7 E.1/E.4; the maps to their finite targets extend to completions.

\paragraph{Finding T04 (source-to-modulus comparison).}
In the recovered LM7\_PUSH E.1/E.4 homology tower, the cusp has order
$12j$ in $G_j=B/(N_B'N_B^j)$. The specifically reported coefficient
checks $j=2,4,8,16,3,9,27,81$, together with the retained $j=6$, therefore
have cusp scalar moduli
\[
        24,48,96,192,36,108,324,972,72.
 \tag{T04.1}
\]
For the correcting characters of LM8\_PUSH C.33--C.38,
\[
 c_{p,k}(u)=\frac{\log(u^{p-1})}{p}\pmod{p^k},
 \qquad p\geq5,\quad k\geq1,
 \tag{T04.2}
\]
the exact scalar conductor is $p^{k+1}$. At $k=1$ this is
$(u^{p-1}-1)/p\bmod p$. The T01 cutoff at that modulus is
$3(k+1)^2(\log p)^2$.

\begin{proof}
E.1 identifies $C=N_B^{\rm ab}\simeq\mathbb Z^{97}$ and records that
$c=[t^{12}]$ is primitive \cite[E.1, proof; E.4]{LM7run}.
The matrix $SR=\left(\begin{smallmatrix}1&-1\\0&1\end{smallmatrix}\right)$
has order twelve even modulo the single sign, since its diagonal
entries cannot become $-1$ modulo twelve. Thus a power $t^r$ can be trivial
in $G_j$ only if $r=12s$. Its value in the kernel $C/jC$ is then
$sc$, zero exactly when $j\mid s$. Hence $\operatorname{ord}_{G_j}(t)=12j$.
This identifies the cusp scalar, not the whole order of $G_j$.
The pure peripheral element $t^2$ has order $6j$, so the half-multiplier
$(\lambda-1)/2$ in the cycle equations is also determined by
$\lambda\bmod12j$.

For (T04.2), if $u/v\in1+p^{k+1}\mathbb Z_p$, then
$\log((u/v)^{p-1})\in p^{k+1}\mathbb Z_p$; division by $p$ makes
the difference zero modulo $p^k$. Thus the character factors through
$U_{p^{k+1}}$. On $u=1+p^k$, expansion of the logarithm gives
\[
 c_{p,k}(1+p^k)\equiv(p-1)p^{k-1}\not\equiv0\pmod{p^k}.
\]
The omitted terms have valuation at least $k$ after division by $p$,
including $k=1$ since $p\geq5$. Therefore it does not factor through
$U_{p^k}$, proving exactness. Reduction of
$\log(1+pa)/p$ modulo $p$ is $a$, which proves the $k=1$ formula.
\end{proof}

More generally, in a finite characteristic free-group target $H$ of
exponent $e$, reduction of $\lambda$ modulo $2e$ determines every
exponentiation by $\lambda$ and by $(\lambda-1)/2$ in the two cycle equations.
When $e$ is odd, reduction modulo $e$ suffices, since two is invertible.
This gives a safe finite scalar parameter once the actual target is
specified; it is not a claim that $2e$ is minimal. For example LM8
C.14 has exponent twelve and explicitly uses multiplier modulus
twenty-four, whereas LM7 D.7 has exponent $101$ in its pure target
and uses $m=0\bmod101$, equivalently $\lambda=1\bmod101$.
The redundant odd lift modulo $202$ gives the same unit group.

The complete recovered-source inventory is
\texttt{GT\_LEVEL\_SCOPE.tsv}. The following are its principal entries;
``frame'' and ``coefficient'' entries are deliberately not promoted
to unspecified scalar quotients.
\begin{center}
\begin{tabular}{p{0.24\linewidth}p{0.23\linewidth}p{0.45\linewidth}}
\hline
Source & Integer or family & Actual meaning and scope\\\hline
LM7 E.1,E.4 & $12j$ & Cusp order in the specified homology quotient.\\
LM7 E.5 & $360$ & Cusp order in $G_6\times A_5$; constructed multiplier one.\\
LM7 E.13,E.14 & $288,576$ & Explicit multiplier-one congruences in the full register.\\
LM7 E.22,E.24 & $72,144$ & Cusp order and the branches $1,73\bmod144$ over retained points.\\
LM7 D.7 & $101$ & Odd-prime multiplier parameter; the tested stratum has multiplier one.\\
LM7 E.50,E.51 & $20,28$ & Actual projective congruence quotients; symplectic tests have multiplier one.\\
LM7 E.52--E.56 & $m,2m,4m$, $(m,6)=1$ & Infinite congruence and exponent families, with specified maps.\\
LM7 E.65 & $12$ & Signed CM frame modulus on $\ker\lambda$, not a cyclotomic character.\\
LM8 C.12 & Every $N$ & Normalized implementer $\left(\begin{smallmatrix}\lambda&t\\0&1\end{smallmatrix}\right)$; determinant is $\lambda$.\\
LM8 C.14 & $24$ & Explicit multiplier modulus for the exponent-twelve $S_4$ envelope.\\
LM8 C.18 & $15$ & Exact domain character $(\lambda/5)(\lambda\bmod3)$.\\
LM8 C.33--C.38 & $p^{k+1}$ & Scalar correcting character (T04.2), not the whole corrected character.\\
LM8 C.49 & $24,60$ & Framed images restricted to multiplier one; no full arbitrary-multiplier image claimed.\\
GT\_TOWER summary & $4,12;\ m=2,3,6;\ 6m=12,18,36$ & Named matrix and power levels only; the full register is unavailable.\\\hline
\end{tabular}
\end{center}

For C.12 the determinant of an arbitrary projective matrix is defined
only modulo squares. It is the displayed unique normalized
representative that retains the full multiplier. For C.49, the orders
$96,48,2880,1440$ of the two framed/arithmetic images are not
cyclotomic moduli. Likewise $T_{24}$ and $T_{48}$ in LM7 E.59/E.62
name targets built from matrix groups of those orders; their subscripts
are not established multiplier moduli. The finite coefficient depths
in the relative nilpotent towers have the same distinction.

\subsection{Direct finite certificates and verification}

\paragraph{Finding T05 (53 exact finite checks).}
For all 53 integers in \texttt{GT\_UNIT\_CERTIFICATES.tsv}, the primes
strictly below $2(\log N)^2$ and coprime to $N$ generate $U_N$.
This assertion is unconditional and finite. A literal word in the
listed prime generators is supplied for every unit in
\texttt{GT\_UNIT\_WORDS.tsv}; there are 4,825 such words in total.
The table also verifies the coefficient-$3$ bound, though no analytic
theorem is needed for these individual computations.

\begin{proof}
The executable \texttt{verify\_gt\_units.py} constructs a subgroup by
successive multiplication and records a nonnegative prime exponent
vector for every reached residue. All arithmetic of group elements is
integer modular arithmetic. The separate read-only
\texttt{check\_unit\_certificates.py} reads the saved tables, without
performing any subgroup search, and re-evaluates each word. It verifies
that each result is a unit, no residue is repeated, and the number
of distinct residues is
$\varphi(N)=N\prod_{p\mid N}(1-1/p)$. It also compares them with
every integer residue coprime to $N$. Thus the displayed products
constitute a certificate of generation, independently of the method
used to discover them. Both executables were run successfully.

Strict bounds are certified with rational intervals, rather than
rounded logarithms. Write $N=2^h r$ with $1\leq r<2$. For rational
$a/b\in[1,2]$, set $t=(a-b)/(a+b)$ and use
\[
 2\sum_{j=0}^{K-1}\frac{t^{2j+1}}{2j+1}
 \ \leq\ \log(a/b)\ \leq\
 2\sum_{j=0}^{K-1}\frac{t^{2j+1}}{2j+1}
       +\frac{2t^{2K+1}}{(2K+1)(1-t^2)}.
 \tag{T05.1}
\]
The error bound follows by replacing every denominator of the
remaining positive terms by $2K+1$, which can only enlarge the sum,
then summing the resulting geometric series. Combining the
intervals for $h\log2$ and $\log r$, and squaring their nonnegative
endpoints, gives rational bounds for $c(\log N)^2$. The program
increases $K$ until both endpoints have the same integer part and
the lower endpoint exceeds that integer. Every proposed prime below
or equal to that integer is therefore strictly below the true bound.
This test terminated for all reported moduli at both $c=2$ and $c=3$.
The printed decimals are only a convenient display of these checks.
\end{proof}

\begin{longtable}{rrrrl}
\caption{Exact finite unit-generation certificates. Decimal bounds are informational; the cutoff column is certified with rational intervals. The lists are generating lists, not claims of minimal cardinality.}\\
\hline
$N$&$\varphi(N)$&$2(\log N)^2$ (display)&Exact cutoff&Prime generators\\\hline
\endfirsthead
\hline
$N$&$\varphi(N)$&$2(\log N)^2$ (display)&Exact cutoff&Prime generators\\\hline
\endhead
2 & 1 & 0.960906 & 0 & $\varnothing$ \\
3 & 2 & 2.413898 & 2 & $2$ \\
4 & 2 & 3.843624 & 3 & $3$ \\
5 & 4 & 5.180581 & 5 & $2$ \\
6 & 2 & 6.420804 & 6 & $5$ \\
7 & 6 & 7.573133 & 7 & $2,3$ \\
8 & 4 & 8.648154 & 8 & $3,5$ \\
9 & 6 & 9.655592 & 9 & $2$ \\
10 & 4 & 10.603796 & 10 & $3$ \\
11 & 10 & 11.499803 & 11 & $2$ \\
12 & 4 & 12.349522 & 12 & $5,7$ \\
13 & 12 & 13.157930 & 13 & $2$ \\
14 & 6 & 13.929247 & 13 & $3$ \\
15 & 8 & 14.667072 & 14 & $2,7$ \\
16 & 8 & 15.374496 & 15 & $3,5$ \\
17 & 16 & 16.054196 & 16 & $2,3$ \\
18 & 6 & 16.708498 & 16 & $5$ \\
19 & 18 & 17.339442 & 17 & $2$ \\
20 & 8 & 17.948824 & 17 & $3,11$ \\
23 & 22 & 19.662648 & 19 & $2,5$ \\
24 & 8 & 20.200052 & 20 & $5,7,13$ \\
25 & 20 & 20.722323 & 20 & $2$ \\
27 & 18 & 21.725081 & 21 & $2$ \\
28 & 12 & 22.207174 & 22 & $3,5$ \\
29 & 28 & 22.677362 & 22 & $2$ \\
30 & 8 & 23.136287 & 23 & $7,11$ \\
31 & 30 & 23.584536 & 23 & $2,3$ \\
36 & 12 & 25.683216 & 25 & $5,7$ \\
48 & 16 & 29.972395 & 29 & $5,7,13$ \\
49 & 42 & 30.292530 & 30 & $2,3$ \\
60 & 16 & 33.527315 & 33 & $7,11,13$ \\
64 & 32 & 34.592617 & 34 & $3,5$ \\
72 & 24 & 36.579746 & 36 & $5,7,13$ \\
81 & 54 & 38.622367 & 38 & $2$ \\
96 & 32 & 41.666549 & 41 & $5,7,13$ \\
101 & 100 & 42.598675 & 42 & $2$ \\
108 & 36 & 43.844706 & 43 & $5,7$ \\
120 & 32 & 45.840154 & 45 & $7,11,13,17$ \\
121 & 110 & 45.999214 & 45 & $2$ \\
144 & 48 & 49.398088 & 49 & $5,7,13$ \\
169 & 156 & 52.631722 & 52 & $2$ \\
192 & 64 & 55.282515 & 55 & $5,7,13$ \\
202 & 100 & 56.355412 & 56 & $3$ \\
288 & 96 & 64.138243 & 64 & $5,7,13$ \\
289 & 272 & 64.216783 & 64 & $2,3$ \\
324 & 108 & 66.833991 & 66 & $5,7$ \\
360 & 96 & 69.292441 & 69 & $7,11,13,17$ \\
361 & 342 & 69.357767 & 69 & $2$ \\
529 & 506 & 78.650592 & 78 & $2,5$ \\
576 & 192 & 80.800209 & 80 & $5,7,13$ \\
841 & 812 & 90.709450 & 90 & $2$ \\
961 & 930 & 94.338145 & 94 & $2,3$ \\
972 & 324 & 94.651073 & 94 & $5,7$ \\
\hline
\end{longtable}


\paragraph{Finding T06 (the remaining evidence boundary).}
The six attached \texttt{SORTED\_*.md} files are summaries, not the
complete row registers their filenames are said to accompany. The
GT\_TOWER\_CHECK summary names a register with 180 rows, but does not
contain that register or the exact definitions of every finite target.
Search and recovery in this run obtained LM8's full 304-row register
and source, and LM7's full 218-row register with only part of its proof
sources. They did not obtain GT\_TOWER\_CHECK's full register or paper.
Consequently T05 proves every listed arithmetic certificate, but is
not described as an exhaustive computation of all finite levels in
the unrecovered 180-row run. The precise remaining gate is to obtain
its finite target definitions and identify their multiplier factors;
the reusable verifier then checks any additional explicit moduli.
For the genuinely infinite families already present in the recovered
sources, T01 applies at every identified modulus, and T03 gives the
stronger unconditional coefficient-$2$ result for the indicated towers.

\begin{proof}
This is a source-coverage finding. The attached summary's own register
count is 180; its displayed statements are truncated and list only
selected rows. The recovered LM7 and LM8 CSV files have 218 and 304
rows respectively. The source inventory records, for each use, whether
the underlying proof, full statement, or only summary was available.
Thus neither a missing row nor an unknown quotient map has been
silently assigned an integer modulus. No analytic assumption can
supply those omitted definitions.
\end{proof}
